% -----------------------------------------------------------------------------
% centriMembSolver — Theory, Implementation and Usage Guide
%
% This document serves as a combined user and developer guide for
% ``centriMembSolver``, an OpenFOAM®-based solver developed for the numerical
% simulation of rotating membrane systems. The documentation is intended to be
% concise, technically precise, and self-contained.
%
% The solver couples incompressible hydrodynamics, solute transport,
% centrifugal acceleration, Coriolis effects, and osmotic phenomena across
% semi-permeable membrane boundaries. The explicit treatment of the Coriolis
% term and its coupling within the PIMPLE iteration are described in detail,
% together with the membrane boundary conditions, numerical algorithms, and
% overall solver architecture.
%
% The manual is intended for advanced users and developers with prior
% experience in OpenFOAM and continuum transport phenomena.
% -----------------------------------------------------------------------------

\documentclass[11pt,a4paper]{article}

% ------------------------- Packages and settings -------------------------
\usepackage[utf8]{inputenc}
\usepackage[T1]{fontenc}
\usepackage{amsmath,amssymb}
\usepackage{siunitx}
\usepackage{graphicx}
\usepackage{booktabs}
\usepackage{tikz}
\usepackage{hyperref}
\usepackage{geometry}
\usepackage{caption}
\usepackage{subcaption}

\geometry{margin=25mm}
\setlength{\parindent}{0pt}
\setlength{\parskip}{1ex plus 0.5ex minus 0.2ex}

\hypersetup{
	colorlinks=true,
	linkcolor=blue,
	urlcolor=blue,
	citecolor=blue
}

% ------------------------ Title and meta information ------------------------
\title{\bfseries centriMembSolver — Theory, Implementation and Usage Guide\\[1ex]
	\large A combined user and developer manual for rotating membrane simulations}

\author{
	V\'itor Geraldes\\
	Department of Chemical Engineering\\
	Instituto Superior Técnico — Universidade de Lisboa\\
	\texttt{vitor.geraldes@tecnico.ulisboa.pt}
}

\date{\today}

\begin{document}
	
	% ======================================================================
	% Page 1 — Title page
	% ======================================================================
	\begin{titlepage}
		\maketitle
		\thispagestyle{empty}
	\end{titlepage}
	
	% ======================================================================
	% Page 2 — License and trademark notice
	% ======================================================================
	\clearpage
	\thispagestyle{plain}
	
	\section*{License and trademarks}

	\textbf{Copyright.} \copyright\ 2026 V\'itor Geraldes.

	\medskip
	\textbf{Software.}
	\texttt{centriMembSolver} is distributed under the GNU General Public
	License version 3 (GPLv3) or later. See the \texttt{LICENSE} file
	distributed with the source code. The software is provided ``as is'',
	WITHOUT WARRANTY OF ANY KIND, express or implied, including but not
	limited to the warranties of merchantability, fitness for a particular
	purpose and noninfringement; see the GPL for details.

	\medskip
	\textbf{Documentation.}
	This manual is licensed under the Creative Commons Attribution 4.0
	International License (CC BY 4.0), unless stated otherwise.

	\medskip
	\textbf{Trademarks.}
	OpenFOAM\textsuperscript{\textregistered} is a registered trademark of
	OpenCFD Limited (ESI Group). Linux\textsuperscript{\textregistered} is
	a registered trademark of Linus Torvalds. ParaView is a trademark of
	Kitware Inc. This work is not affiliated with, endorsed by, or
	sponsored by OpenCFD Limited, the Linux Foundation, or Kitware Inc.

	% ======================================================================
	% Page 3 — How to cite
	% ======================================================================
	\clearpage
	\thispagestyle{plain}

	\section*{How to cite}
	\addcontentsline{toc}{section}{How to cite}

	\noindent
	\texttt{centriMembSolver} is distributed as an open-source platform for
	the community of rotating-membrane simulations.  Users who publish work
	based on the solver are encouraged to cite the versioned software
	release directly:

	\begin{quote}
		Geraldes, V., \emph{centriMembSolver: a finite-volume solver for
		rotating membrane systems built on OpenFOAM\textsuperscript{\textregistered}
		technology}, version 1.1.0, Zenodo,
		\textbf{[DOI to be inserted upon Zenodo deposit]}, 2026.
	\end{quote}

	A BibTeX entry is provided in the source repository
	(\href{https://github.com/vitormgeraldes/centriMembSolver}{\texttt{github.com/vitormgeraldes/centriMembSolver}}).

	\medskip
	\noindent
	\textbf{Application example.}  A first quantitative application of
	\texttt{centriMembSolver} to the prediction of concentration
	polarisation in Ortho-Centrifugal Membrane Filters is described in
	Geraldes \emph{et al.} (\emph{Predicting concentration polarization
	in Ortho-Centrifugal Membrane Filters for NF/RO membrane screening},
	submitted to \emph{Journal of Membrane Science}).  That work
	documents a buoyancy-driven mass-transfer correlation derived from a
	direct numerical simulation dataset produced with this solver and is
	the appropriate reference for users adopting the same physical
	framework.

	\medskip
	\noindent
	\textbf{Methodological extensions.}  Users developing new boundary
	conditions, transport closures, multi-physics couplings, or new
	application studies are encouraged to publish their own descriptions
	and to reference the \texttt{centriMembSolver} software DOI as the
	implementation platform.  Pull requests with new tutorials, BC
	classes or transport models are welcome via the source repository.

	% ======================================================================
	% Page 4 — Abstract
	% ======================================================================
	\clearpage
	\thispagestyle{plain}

	\section*{Abstract}
	\addcontentsline{toc}{section}{Abstract}
	
	\noindent
	\texttt{centriMembSolver} is a finite–volume solver for coupled solvent and
	solute transport through rotating membranes. It extends
	\texttt{buoyantBoussinesqPimpleFoam} to include solute transport and semi-permeable membrane boundary conditions,
	concentration-dependent viscosity, and an explicit treatment of the
	Coriolis force consistent with the Boussinesq buoyancy formulation under
	centrifugal acceleration. This guide documents the underlying theory,
	the implementation of membrane flux models, and the overall solver
	architecture. A worked example, \emph{rotatingSlit}, is included to
	illustrate case setup and post-processing. The document is intended both
	for end-users running membrane simulations and for developers wishing to
	extend the solver to new flux laws or multi-physics couplings.
	
	% ======================================================================
	% Page 4 — Table of contents
	% ======================================================================
	\clearpage
	\thispagestyle{plain}
	\tableofcontents

	% ======================================================================
	% Notation
	% ======================================================================
	\clearpage
	\thispagestyle{plain}
	\section*{Notation}
	\addcontentsline{toc}{section}{Notation}

	\noindent
	The following table maps the mathematical symbols used throughout this
	manual to the corresponding OpenFOAM field names and dictionary keys.

	\renewcommand{\arraystretch}{1.15}
	\begin{tabular}{@{}lp{7cm}ll@{}}
		\toprule
		Symbol & Quantity & Field / dict key & Units\\
		\midrule
		$\boldsymbol{U}$ & velocity field (rotating frame) & \texttt{U} & m\,s$^{-1}$\\
		$p_\mathrm{rgh}$ & modified (dynamic) pressure & \texttt{p\_rgh} & Pa\\
		$\mathrm{gh}$ & combined gravitational + centrifugal hydrostatic potential per unit mass & \texttt{gh} & m$^2$\,s$^{-2}$\\
		$\rho_k$ & dimensionless density ratio $\rho/\rho_0$ & \texttt{rhok} & --\\
		$\mathbf{r}$ & position vector relative to the rotation axis & --- & m\\
		$\mathbf{r}_{\!\perp}$ & component of $\mathbf{r}$ perpendicular to the rotation axis & --- & m\\
		$\boldsymbol{\Omega}$ & angular velocity vector (along rotation axis) & \texttt{SRFProperties} (\texttt{rpm}, \texttt{axis}) & rad\,s$^{-1}$\\
		$\Omega$ & angular speed, $\Omega = \lvert\boldsymbol{\Omega}\rvert$ & --- & rad\,s$^{-1}$\\
		$r_m$ & membrane orbit radius (for centrifugal NF/RO cartridges) & --- & m\\
		$a_m$ & centrifugal acceleration at the membrane, $a_m = \Omega^2 r_m$ & --- & m\,s$^{-2}$\\
		\midrule
		$C_A$ & primary solute concentration & \texttt{CA} & kg\,m$^{-3}$\\
		$D_A(C_A)$ & primary solute diffusivity (concentration--dependent) & \texttt{speciesData} (\texttt{DAB0..4}) & m$^2$\,s$^{-1}$\\
		$\nu(C_A)$ & kinematic viscosity (concentration--dependent) & \texttt{speciesData} (\texttt{nu0..2}) & m$^2$\,s$^{-1}$\\
		$\beta$ & solutal contraction coefficient $(1/\rho_0)(\partial\rho/\partial C_A)$ & \texttt{speciesData} (\texttt{beta}) & (kg\,m$^{-3}$)$^{-1}$\\
		$\rho_0$ & reference density (water at $T_\mathrm{ref}$) & \texttt{speciesData} (\texttt{rho0}) & kg\,m$^{-3}$\\
		$\pi(C_A)$ & osmotic pressure & via \texttt{virialCoeffs} & Pa\\
		\midrule
		$C_T$ & tracer (passive secondary solute) concentration & \texttt{CT} & kg\,m$^{-3}$\\
		$D_T$ & tracer diffusivity (constant) & \texttt{transportProperties} (\texttt{DT0}) & m$^2$\,s$^{-1}$\\
		\midrule
		$J_v$ & solvent (water) flux through membrane & patch field \texttt{Jv} & m\,s$^{-1}$\\
		$J_s$ & solute mass flux through membrane & patch field \texttt{Js} & kg\,(m$^2$\,s)$^{-1}$\\
		$J_{s,T}$ & tracer mass flux through membrane & patch field \texttt{JsT} & kg\,(m$^2$\,s)$^{-1}$\\
		$C_w$, $C_p$ & solute concentrations at the membrane wall (feed side) and permeate side & internal $/$ \texttt{CAp} & kg\,m$^{-3}$\\
		$C_f$, $C_b$ & feed-side and bulk reference concentrations & dict (\texttt{CAb}) & kg\,m$^{-3}$\\
		\midrule
		$A_h$ & hydraulic permeability of membrane & \texttt{Ah} & m\,(Pa\,s)$^{-1}$\\
		$\sigma$ & solute reflection coefficient (Spiegler--Kedem) & \texttt{sigma} & --\\
		$\sigma_T$ & tracer reflection coefficient (default 1) & \texttt{sigmaT} & --\\
		$B$ & solute permeability coefficient (Kedem--Katchalsky) & \texttt{B} & m\,s$^{-1}$\\
		$B_T$ & tracer permeability coefficient & \texttt{BT} & m\,s$^{-1}$\\
		$R_\mathrm{int}$, $R_\mathrm{obs}$ & intrinsic and observed solute rejection & \texttt{Rint}, \texttt{Robs} & --\\
		$R_{T,\mathrm{int}}$, $R_{T,\mathrm{obs}}$ & intrinsic and observed tracer rejection & \texttt{RTint}, \texttt{RTobs} & --\\
		$\Gamma$ & solute concentration polarisation $(C_w - C_b)/C_b$ & volume field \texttt{Gama} & --\\
		$\Gamma_T$ & tracer concentration polarisation & volume field \texttt{GamaT} & --\\
		\bottomrule
	\end{tabular}
	\renewcommand{\arraystretch}{1.0}

	% ======================================================================
	% Page 5+ — Main content
	% ======================================================================
	\clearpage

	% ======================================================================

\section{Introduction}

Rotating membrane systems combine fluid flow, solute transport and
selective separation in a centrifugal field.  Applications range from
ultra–scale–down (USD) process development, where small membrane coupons
are rotated at high speed to emulate industrial centrifugation, to
high-throughput screening of nanofiltration and reverse osmosis (NF/RO)
membranes in centrifugal filter cartridges.  In these systems a
pressure gradient drives solvent through a semi–permeable barrier while
a concentration difference induces osmotic counterflow.  Rotation adds
centrifugal and Coriolis effects that modify the hydrodynamics and
solute distribution.

\texttt{centriMembSolver} was developed to provide an open and
customisable platform for simulating such coupled phenomena.  Built on
OpenFOAM® v2506, the solver solves the incompressible Navier–Stokes
equations with source terms for rotation, a scalar transport equation
for solute concentration and custom patch fields to model membrane
fluxes.  The Coriolis acceleration is treated as an explicit source
term in the momentum equation, refreshed at every PIMPLE outer
iteration; this is sufficient under the small CFL-bounded time steps
characteristic of rotating membrane DNS.  In addition, the membrane
boundary conditions
implement well established flux laws—Spiegler–Kedem, constant flux and
target average flux for the solvent, and intrinsic rejection, solute
permeability or observed rejection for the solute.  Diagnostic
quantities such as observed rejection and concentration polarisation are
computed automatically. As of the 2026 release, the solver has been extended to simulate a primary solute A and a passive tracer T with separate diffusivities and membrane flux models, enabling tracer-based validation experiments.

The following sections present the governing equations (including
viscous and osmotic effects), derive the robust Coriolis algorithm,
detail the boundary conditions, describe the solver structure and
illustrate usage through a case study.  Developer notes at the end
provide guidance on extending the code and best practice for stable
simulations.

An example application of the solver, including validation against
experimental data, is summarised on page~\pageref{sec:cases} and
described in detail in the work cited under \emph{How to cite}.

\subsection*{Quick start}
\addcontentsline{toc}{subsection}{Quick start}

For users who want to run a case immediately, the following commands
compile the libraries and the solver and execute the bundled
two-dimensional tutorial \texttt{rotatingSlit}.  \texttt{<repo>}
denotes the root of the cloned source tree, and an active OpenFOAM
environment (\texttt{source} of \texttt{etc/bashrc}) is assumed.

\paragraph{Serial (single-process) workflow.}
\begin{enumerate}
	\item Build the membrane boundary-condition and transport
	libraries, then the solver:
	\begin{verbatim}
cd <repo>/src     && wmake all libso
cd <repo>/solver  && wmake
	\end{verbatim}

	\item Run the bundled tutorial:
	\begin{verbatim}
cd <repo>/tutorials/rotatingSlit
blockMesh
centriMembSolver
	\end{verbatim}

	\item Visualise the results in ParaView (requires the
	\texttt{paraFoam} reader shipped with OpenFOAM):
	\begin{verbatim}
paraFoam
	\end{verbatim}
\end{enumerate}

\paragraph{Parallel (MPI) workflow.}
For three-dimensional cases the solver runs under MPI in the standard
OpenFOAM idiom:
\begin{verbatim}
decomposePar                                     # split mesh+fields
mpirun -np <N> centriMembSolver -parallel        # run on N processes
reconstructPar                                   # merge fields for paraFoam
\end{verbatim}
The number of processes \texttt{<N>} is set by
\texttt{numberOfSubdomains} in
\texttt{system/decomposeParDict}; standard \texttt{simple} or
\texttt{scotch} decomposition methods are supported.

Section~\ref{sec:cases} walks through the bundled tutorial file by
file and explains the diagnostic outputs.

% ===========================================================================


\section{Theory and Governing Equations}

\subsection{Hydrodynamics}

In a non-inertial reference frame rotating with angular velocity
\(\boldsymbol{\Omega}\), the incompressible Navier–Stokes equations read
\begin{equation}
	\nabla\cdot\mathbf{U} = 0,\qquad
	\rho \frac{\partial\mathbf{U}}{\partial t} + \rho\,\nabla\cdot(\mathbf{U}\mathbf{U})
	= -\nabla p + \nabla\cdot(\mu_{\mathrm{eff}}\nabla\mathbf{U})
	+ \rho\,\mathbf{f}_{\mathrm{rot}},
\end{equation}
where \(\mathbf{U}\) is the velocity, \(p\) the kinematic pressure
(\(p=p_{\mathrm{phys}}/\rho\)), and the effective viscosity
\(\mu_{\mathrm{eff}} = \rho(\nu + \nu_t)\) includes molecular and—if
selected—turbulent viscosity.

The rotational body force \(\mathbf{f}_{\mathrm{rot}}\) contains the
two pseudo-forces appearing in a uniformly rotating frame, the
centrifugal and Coriolis terms,
\begin{equation}
	\mathbf{f}_{\mathrm{rot}}
	= -\,\boldsymbol{\Omega}\times(\boldsymbol{\Omega}\times\mathbf{r})
	- 2\,\boldsymbol{\Omega}\times\mathbf{U},
\end{equation}
with \(\mathbf{r}\) the position vector relative to the rotation axis.
The first term is conservative and points radially outward,
\(-\boldsymbol{\Omega}\times(\boldsymbol{\Omega}\times\mathbf{r})
= +\Omega^2\,\mathbf{r}_{\!\perp}\); the second is non-conservative.
Pressure--velocity coupling is handled by the PIMPLE algorithm
(Pressure-Implicit Method for Pressure-Linked Equations, the standard
OpenFOAM transient pressure-correction scheme), as implemented in
\texttt{UEqn.H} and \texttt{pEqn.H}.







\subsubsection{Modified pressure, hydrostatic potential and Boussinesq buoyancy}

Rotating membrane devices exhibit very large hydrostatic pressure
gradients, dominated by centrifugal effects. The physical pressure
\(p\) may increase by several bars over a few centimetres of radius,
while the dynamic variations induced by fluid motion are much smaller.
If \(p\) is solved directly, the linear systems become ill–conditioned
and the convergence of the PIMPLE algorithm deteriorates.

To avoid this, \texttt{centriMembSolver} follows the standard OpenFOAM
strategy and works with three auxiliary fields:

\begin{itemize}
	\item \texttt{p\_rgh}: modified (dynamic) pressure
	\item \texttt{gh}: hydrostatic potential per unit mass
	\item \texttt{rhok}: dimensionless density ratio
\end{itemize}

The solver only solves for the dynamic component \texttt{p\_rgh}; the
hydrostatic contribution is treated analytically via \texttt{gh} and
\texttt{rhok}.

\paragraph{Boussinesq density field}

Under the Boussinesq approximation, the density depends weakly on
solute concentration,
\begin{equation}
	\rho(\boldsymbol{x}) = \rho_0\,\rho_k(\boldsymbol{x}),
	\qquad
	|\rho_k - 1| \ll 1,
\end{equation}
where the code field \texttt{rhok} corresponds to the scalar
\(\rho_k\). The Boussinesq approximation replaces \(\rho\) by the
constant reference value \(\rho_0\) in every term of the momentum
equation \emph{except} the body--force (buoyancy) term, where the
actual density \(\rho = \rho_0\rho_k\) is retained.  This linearises
the inertia, viscous and Coriolis terms while capturing the
leading-order density-driven body force.  Consistently with this
approximation, the velocity field is treated as divergence--free,
\begin{equation}
	\nabla\cdot\boldsymbol{U} = 0,
	\qquad
	\frac{\partial \rho}{\partial t} = 0,
\end{equation}
but density variations (\(\rho_k \neq 1\)) still generate buoyancy forces.

\paragraph{Effective acceleration in the rotating frame}

Let \(\boldsymbol{x}\) denote the position vector in the rotating frame
and \(\boldsymbol{x}_0\) a point on the rotation axis. Both are read
from \texttt{constant/SRFProperties} (Single Reference Frame: the
OpenFOAM dictionary that specifies the rotation axis through
\texttt{origin} and \texttt{axis}, and the rotation rate through
\texttt{rpm}). A fluid particle at rest in this frame is
subject to gravity and centrifugal acceleration:
\begin{equation}
	\boldsymbol{a}_{\mathrm{eff}}(\boldsymbol{x})
	=
	\boldsymbol{g}
	-
	\boldsymbol{\Omega}\times
	\bigl(
	\boldsymbol{\Omega}\times(\boldsymbol{x}-\boldsymbol{x}_0)
	\bigr),
	\label{eq:a_eff}
\end{equation}
where \(\boldsymbol{g}\) is the gravitational acceleration and
\(\boldsymbol{\Omega}\) the angular velocity vector.  Note that
\(-\boldsymbol{\Omega}\times(\boldsymbol{\Omega}\times\mathbf{r}_\perp)
= +\Omega^2\,\mathbf{r}_\perp\), where \(\mathbf{r}_\perp\) is the
component of \(\mathbf{r}\) perpendicular to the rotation axis: the
centrifugal acceleration points radially outward.

The Coriolis term \(-2\,\boldsymbol{\Omega}\times\boldsymbol{U}\) is
non–conservative and is therefore \emph{not} included in
\(\boldsymbol{a}_{\mathrm{eff}}\).

\paragraph{Hydrostatic potential \texttt{gh}}

The key point is that the effective acceleration \(\boldsymbol{a}_{\rm eff}\)
in \eqref{eq:a_eff} is conservative: it can be written as the gradient
of a scalar potential per unit mass. OpenFOAM stores this potential in
the field \texttt{gh}, defined so that
\begin{equation}
	\nabla (\mathrm{gh}) = \boldsymbol{a}_{\mathrm{eff}}.
\end{equation}

A convenient explicit expression is
\begin{equation}
	\mathrm{gh}(\boldsymbol{x})
	=
	\boldsymbol{g}\cdot(\boldsymbol{x}-\boldsymbol{x}_{\mathrm{ref}})
	+ \frac{1}{2}
	\left|
	\boldsymbol{\Omega}\times(\boldsymbol{x}-\boldsymbol{x}_0)
	\right|^2
	- \mathrm{gh}_{\mathrm{ref}},
	\label{eq:gh_def}
\end{equation}
where \(\boldsymbol{x}_{\mathrm{ref}}\) is a reference point and
\(\mathrm{gh}_{\mathrm{ref}}\) is a constant offset. This offset does not
affect the physics, since only gradients of \(\mathrm{gh}\) enter the
equations.

\medskip
\noindent
\textbf{Physical meaning.}
\(\mathrm{gh}\) is the hydrostatic potential per unit mass, combining
gravitational and centrifugal contributions. Surfaces of constant
\(\mathrm{gh}\) correspond to equilibrium surfaces of a fluid at rest
under gravity and rotation.

\medskip
\noindent
\textbf{Why store the centrifugal potential in \texttt{gh}.}
Both gravity and the centrifugal acceleration of the rotating frame are
conservative body forces and can therefore be absorbed into a single
hydrostatic potential, with their density-fluctuation counterparts
retained as the buoyancy drivers. This is the classical construction
for rotating Boussinesq flows (Chandrasekhar 1961; Lopez, Marques \&
Avila 2013): explicitly,
\(-\rho_0\,\boldsymbol{\Omega}\times(\boldsymbol{\Omega}\times\mathbf{r})
= \nabla\!\bigl(\tfrac12\rho_0|\boldsymbol{\Omega}\times\mathbf{r}|^2\bigr)\)
is absorbed into the pressure gradient, leaving
\(-\rho'\,\boldsymbol{\Omega}\times(\boldsymbol{\Omega}\times\mathbf{r})\)
to drive centrifugal buoyancy. In \texttt{centriMembSolver} we realise
this within OpenFOAM's existing \texttt{gh}/\texttt{p\_rgh}
infrastructure (Eq.~\eqref{eq:gh_def}); only the non-conservative
Coriolis term remains as an explicit source in the momentum equation.
The practical motivation is numerical conditioning: at high rotation
rates the centrifugal hydrostatic pressure can dominate the absolute
pressure field by orders of magnitude (\(\sim 60\) bar in
high-throughput screening cartridges), and absorbing it into
\texttt{p\_rgh} keeps the pressure–velocity coupling well conditioned.

\paragraph{Modified pressure \(p_{\mathrm{rgh}}\)}

In the solver, the physical pressure is reconstructed from the three
fields \texttt{p\_rgh}, \texttt{gh} and \texttt{rhok} as
\begin{equation}
	p(\boldsymbol{x})
	=
	\rho_0\bigl(
	p_{\mathrm{rgh}}(\boldsymbol{x})
	+ \rho_k(\boldsymbol{x})\,\mathrm{gh}(\boldsymbol{x})
	\bigr)
	+ p_{\mathrm{ref}},
	\label{eq:p_reconstruct}
\end{equation}
where \(p_{\mathrm{rgh}}\) corresponds to the OpenFOAM field
\texttt{p\_rgh} and \(p_{\mathrm{ref}}\) is a constant reference pressure.

Neglecting the constant \(p_{\mathrm{ref}}/\rho_0\), this can be
rearranged as
\begin{equation}
	p_{\mathrm{rgh}}
	=
	\frac{p}{\rho_0}
	- \rho_k\,\mathrm{gh}.
	\label{eq:p_rgh_def}
\end{equation}

Thus \(p_{\mathrm{rgh}}\) represents the \emph{dynamic} part of the
pressure: the large hydrostatic background is removed and only the
smaller fluctuations are handled by the linear solver.

\paragraph{Momentum equation in \(p_{\mathrm{rgh}}\) form}

Applying the Boussinesq approximation as stated above (\(\rho\to\rho_0\)
in inertia, viscous and Coriolis terms; \(\rho=\rho_0\rho_k\) retained
only in the body--force term), the incompressible rotating--frame
Navier--Stokes equation becomes
\begin{equation}
	\rho_0
	\left(
	\frac{\partial \boldsymbol{U}}{\partial t}
	+ \boldsymbol{U}\cdot\nabla\boldsymbol{U}
	+ 2\,\boldsymbol{\Omega}\times\boldsymbol{U}
	\right)
	=
	-\nabla p
	+ \nabla\cdot\bigl(\mu_{\mathrm{eff}}\nabla\boldsymbol{U}\bigr)
	+ \rho_0\rho_k\,\boldsymbol{a}_{\mathrm{eff}}.
	\label{eq:NS_raw}
\end{equation}

\paragraph{Step-by-step transformation to the \(p_{\mathrm{rgh}}\)-based form}

Using the definition \eqref{eq:p_rgh_def},
\[
p
= \rho_0\left(p_{\mathrm{rgh}} + \rho_k\,\mathrm{gh}\right),
\]
the pressure gradient becomes
\begin{align}
	-\nabla p
	&= -\rho_0 \nabla\!\left(p_{\mathrm{rgh}} + \rho_k\,\mathrm{gh}\right)
	\nonumber\\
	&= -\rho_0\nabla p_{\mathrm{rgh}}
	-\rho_0\nabla(\rho_k\,\mathrm{gh}).
\end{align}
Applying the product rule,
\begin{equation}
	\nabla(\rho_k\,\mathrm{gh})
	=
	\mathrm{gh}\,\nabla\rho_k
	+
	\rho_k\,\nabla(\mathrm{gh})
	=
	\mathrm{gh}\,\nabla\rho_k
	+
	\rho_k\,\boldsymbol{a}_{\mathrm{eff}},
\end{equation}
where we used \(\nabla(\mathrm{gh}) = \boldsymbol{a}_{\mathrm{eff}}\).
Thus
\begin{equation}
	-\nabla p
	=
	-\rho_0\nabla p_{\mathrm{rgh}}
	-\rho_0\,\mathrm{gh}\,\nabla\rho_k
	-\rho_0\rho_k\,\boldsymbol{a}_{\mathrm{eff}}.
	\label{eq:pressure_expanded}
\end{equation}

Substituting \eqref{eq:pressure_expanded} into \eqref{eq:NS_raw} and
cancelling the \(\rho_0\rho_k\,\boldsymbol{a}_{\mathrm{eff}}\) terms on
both sides yields
\begin{equation}
	\rho_0
	\left(
	\frac{\partial \boldsymbol{U}}{\partial t}
	+ \boldsymbol{U}\cdot\nabla\boldsymbol{U}
	+ 2\,\boldsymbol{\Omega}\times\boldsymbol{U}
	\right)
	=
	- \rho_0\nabla p_{\mathrm{rgh}}
	- \rho_0\,\mathrm{gh}\,\nabla\rho_k
	+ \nabla\cdot\bigl(\mu_{\mathrm{eff}}\nabla\boldsymbol{U}\bigr).
	\label{eq:NS_prgh_rho0}
\end{equation}

Since \(\rho_0\) is constant, we can divide the entire equation by
\(\rho_0\) and introduce the effective kinematic viscosity
\(\nu_{\mathrm{eff}} = \mu_{\mathrm{eff}}/\rho_0\) to obtain the
canonical form used in the solver:
\begin{equation}
	\frac{\partial \boldsymbol{U}}{\partial t}
	+ \boldsymbol{U}\cdot\nabla\boldsymbol{U}
	+ 2\,\boldsymbol{\Omega}\times\boldsymbol{U}
	=
	- \nabla p_{\mathrm{rgh}}
	- \mathrm{gh}\,\nabla\rho_k
	+ \nabla\cdot\bigl(\nu_{\mathrm{eff}}\nabla\boldsymbol{U}\bigr).
	\label{eq:NS_prgh}
\end{equation}
Note that \(\nu_{\mathrm{eff}} = \nu(C_A) + \nu_t\) is retained
\emph{inside} the divergence operator because the molecular kinematic
viscosity \(\nu(C_A)\) is concentration--dependent (Eq.~(8)).

The hydrostatic load has disappeared from the unknown pressure
(\(p_{\mathrm{rgh}}\)), and only the term proportional to
\(\nabla\rho_k\) remains as a Boussinesq buoyancy force.

\paragraph{Discrete buoyancy term in \texttt{centriMembSolver}}

The continuous term
\(-\,\mathrm{gh}\,\nabla\rho_k\)
is discretised in OpenFOAM using face–based gradients:
\begin{equation}
	-\,\mathrm{gh}\,\nabla\rho_k
	\quad\Longrightarrow\quad
	-\,\mathrm{gh}_f\,\mathrm{snGrad}(\rho_k),
\end{equation}
where \(\mathrm{gh}_f\) is the potential field interpolated to the
faces, and \(\mathrm{snGrad}(\rho_k)\) is the surface–normal gradient of
\(\rho_k\). In the actual code this appears as
\begin{verbatim}
	- ghf * fvc::snGrad(rhok)
\end{verbatim}
inside the reconstruction of the right–hand side of the momentum
equation.

\paragraph{Summary and physical interpretation}

\begin{itemize}
	\item \(\mathrm{gh}\) is the hydrostatic potential per unit mass:
	it combines gravitational and centrifugal contributions and is
	computed from geometry, \(\boldsymbol{g}\) and \(\boldsymbol{\Omega}\).
	\item \(\rho_k\) (field \texttt{rhok}) encodes density variations
	relative to \(\rho_0\) and is responsible for Boussinesq
	buoyancy.
	\item \(p_{\mathrm{rgh}}\) (field \texttt{p\_rgh}) is the pressure
	variable actually solved by PIMPLE; its gradients are much
	smaller than those of the physical pressure \(p\).
	\item The buoyancy force appears as
	\(-\,\mathrm{gh}\,\nabla\rho_k\), discretised in the code
	as \texttt{- ghf*fvc::snGrad(rhok)}.
	\item The physical pressure used for membrane transport and
	post–processing is reconstructed via \eqref{eq:p_reconstruct}.
\end{itemize}

\subsubsection{Turbulence modelling}

Although centrifugal membrane filtration can operate in laminar,
transitional or weakly turbulent regimes, \texttt{centriMembSolver}
allows the user to select any turbulence model available in the
\texttt{RASModel} hierarchy of OpenFOAM.  The model is instantiated via
the usual run-time selection mechanism in \texttt{createFields.H}:
\begin{equation}
	\nu_{\mathrm{eff}}
	= \nu + \nu_t(\mathbf{U}, k, \varepsilon, \dots).
\end{equation}

Typical configurations rely on:
\begin{itemize}
	\item direct laminar modelling (\texttt{simulationType laminar;})
	\item one–equation eddy–viscosity models such as Spalart–Allmaras
	\item two–equation models such as standard or low-Re \(k{-}\varepsilon\)
	~and \(k{-}\omega\)
\end{itemize}
with parameters defined in \texttt{constant/turbulenceProperties}.
Turbulent mass transport uses an eddy diffusivity
\(D_t = \nu_t / \mathrm{Sc}_t\), where \(\mathrm{Sc}_t\) is the turbulent
Schmidt number.

\subsubsection{Concentration-dependent viscosity}

Many aqueous solute solutions (salts, sugars, cryoprotectants,
cryoprotectant carriers) exhibit a strong dependence of viscosity on
the solute concentration \(C_A\).  \texttt{centriMembSolver}
incorporates a polynomial parameterisation:
\begin{equation}
	\nu(C_A) = \nu_0 + \nu_1 C_A + \nu_2 C_A^2 + \cdots,
\end{equation}
with coefficients \(\nu_0,\nu_1,\nu_2\) set in
\texttt{transportProperties}.  A clipping operator prevents unphysical
values at high concentrations.



\subsection{Solute and tracer transport}

In the current release, the solver solves two advective–diffusive transport equations: one for a primary solute concentration \(C_A\) and one for a passive tracer concentration \(C_T\). Both species are advected by the same velocity field \(\mathbf{U}\) but may have different diffusion coefficients.  The governing equations read
\begin{align}
  \frac{\partial C_A}{\partial t} + \nabla\cdot(C_A\mathbf{U}) &= \nabla\cdot(D_{A,\mathrm{eff}}\,\nabla C_A),\\[0.5ex]
  \frac{\partial C_T}{\partial t} + \nabla\cdot(C_T\mathbf{U}) &= \nabla\cdot(D_{T,\mathrm{eff}}\,\nabla C_T),
\end{align}
where \(D_{A,\mathrm{eff}} = D_{AB}(C_A) + D_t\) is the effective diffusivity of the solute including the concentration–dependent molecular diffusivity \(D_{AB}\) and a turbulent contribution \(D_t\) (if enabled), and \(D_{T,\mathrm{eff}} = D_{T,0} + D_t\) is the effective diffusivity of the tracer with \(D_{T,0}\) constant and independent of concentration. Membrane boundary conditions for the solute and tracer are described in \S\ref{sec:bc-solute} and \S\ref{sec:bc-tracer}, respectively.

\subsection{Coriolis treatment in PIMPLE}
\label{sec:coriolis}

The Coriolis pseudo-force \(\mathbf{f}_{C} = -2\,\boldsymbol{\Omega}\times\mathbf{U}\)
is non-conservative and is therefore added explicitly to the momentum
equation as a source term, evaluated with the velocity available at
the current PIMPLE outer iteration
\(\mathbf{U}^{(k)}\):
\begin{equation}
	\mathbf{f}_{C}^{\,(k)} = -2\,\boldsymbol{\Omega}\times\mathbf{U}^{(k)} .
	\label{eq:coriolis-explicit}
\end{equation}
The explicit treatment is appropriate at the small time steps imposed
by the convective CFL constraint of DNS in rotating membrane systems
and avoids the additional matrix structure that an implicit
discretisation would require. For DNS runs the
\texttt{momentumPredictor} is disabled to avoid spurious Courant
excursions in strongly rotating flows; a single outer corrector per
time step is normally sufficient under the CFL-bounded \(\Delta t\).
The corresponding code-level snippet is shown together with the rest
of the momentum loop in \S\ref{sec:solver-structure}.

\section{Membrane Boundary Conditions}

Membrane patches couple the feed region to the permeate side by
imposing fluxes of solvent and solute according to physical models.
\texttt{centriMembSolver} implements two specialised patch field
classes: \texttt{membraneSolventFluxFvPatchVectorField} for the
velocity and \texttt{membraneSoluteFluxFvPatchScalarField} for the
concentration.  The feed side of the membrane is the patch
corresponding to the field values; permeate side conditions are
either fixed (for constant pressure/concentration) or inferred from
the flux law.

\subsection{Membrane physics}
\label{sec:bc-theory}

This subsection collects all the membrane-transport physics that
underlies the implementations of \S\ref{sec:bc-solvent}--\S\ref{sec:bc-tracer}.
It can be skipped on a first reading by users who only want dictionary
syntax; the implementation subsections forward-reference equations
defined here when needed.

\paragraph{Driving forces and osmotic pressure.}
Transmembrane transport is driven by hydrostatic
\(\Delta p = p_f - p_p\) and osmotic
\(\Delta\pi = \pi(C_f) - \pi(C_p)\) differences. The osmotic pressure
is represented by a virial expansion in solute concentration,
\begin{equation}
	\pi(C_A) = \sum_{k=1}^{N} a_k\,C_A^k ,
	\label{eq:osmotic-virial}
\end{equation}
with mass-based coefficients \(a_k\) of units
\si{Pa.m^{3k}.kg^{-k}} (\(C_A\) in \si{kg.m^{-3}}). At infinite
dilution the linear coefficient reduces to the van~'t~Hoff value
\(a_1 = i\,RT/M\), where \(i\) is the dissociation factor, \(R\) the
gas constant, \(T\) absolute temperature, and \(M\) the solute molar
mass; the higher-order \(a_k\) capture deviations from ideality.

\paragraph{Spiegler--Kedem solvent flux.}
The solvent (water) flux \(J_v\) [m/s] follows
\begin{equation}
	J_v = A_h\bigl(\Delta p - \sigma\,\Delta\pi\bigr),
	\label{eq:jv-s-k}
\end{equation}
with hydraulic permeability \(A_h\) (dictionary key \texttt{Ah}) and
reflection coefficient \(\sigma\) (dictionary key \texttt{sigma}).

\paragraph{Kedem--Katchalsky solute flux.}
The solute mass flux \(J_s\) [kg\,(m\(^2\)\,s)\(^{-1}\)] is given by
\begin{equation}
	J_s = (1 - \sigma)\,J_v\,C_w + B\,(C_w - C_p),
	\label{eq:js-s-k}
\end{equation}
where \(B\) is the intrinsic solute permeability (dictionary key
\texttt{B}), \(C_w\) the wall concentration on the feed side, and
\(C_p\) the permeate concentration. The intrinsic rejection
coefficient is, by definition,
\begin{equation}
	R_\mathrm{int} \equiv 1 - C_p/C_w .
	\label{eq:Rint-def}
\end{equation}

\paragraph{Local mass balance at the membrane face.}
The membrane patch separates a feed-side bulk fluid (resolved by the
DNS) from a permeate side. At the wall, the convective influx of
solute, the diffusive removal back into the bulk feed, and the
permeate outflux must balance:
\begin{equation}
	J_v\,C_w \;-\; D_A\,\frac{\partial C_A}{\partial n}\bigg|_{w}
	\;=\; J_s \;=\; J_v\,C_p ,
	\label{eq:bc-balance}
\end{equation}
where \(\partial/\partial n\) is the outward normal derivative on the
feed-side patch (\texttt{snGrad} in OpenFOAM, defined as
\((C_w - C_i)/\delta_n\) for the cell adjacent to the face, with
\(\delta_n\) the cell-centre--to--face distance). Equation
\eqref{eq:bc-balance} couples \(C_w\) to the resolved interior
field; it makes \emph{no} use of a phenomenological mass-transfer
coefficient because the boundary layer is resolved by the DNS.

The \emph{closure} required to solve \eqref{eq:bc-balance} for
\(C_w\) is provided differently by the two membrane-physics
specialisations:
\begin{itemize}
	\item \textbf{Intrinsic-rejection closure.}
	The membrane is assumed to reject a fixed fraction
	\(R_\mathrm{int}\) of the solute regardless of operating point,
	\begin{equation}
		C_p = (1 - R_\mathrm{int})\,C_w .
		\label{eq:Cp-from-Rint}
	\end{equation}
	Substituting \eqref{eq:Cp-from-Rint} into \eqref{eq:bc-balance}
	gives the Robin-type relation
	\begin{equation}
		D_A\,\frac{\partial C_A}{\partial n}\bigg|_{w}
		= R_\mathrm{int}\,J_v\,C_w .
		\label{eq:Rint-Robin}
	\end{equation}

	\item \textbf{Solute-permeability (Spiegler--Kedem) closure.}
	Combining the SK solute flux \eqref{eq:js-s-k} with
	\(J_s = J_v\,C_p\) and solving for \(C_p\) gives
	\begin{equation}
		C_p = \frac{B + (1-\sigma)\,J_v}{B + J_v}\,C_w
		    \;\equiv\; \alpha(J_v;\sigma,B)\,C_w ,
		\label{eq:Cp-from-SK}
	\end{equation}
	so that the effective intrinsic rejection of an SK membrane is
	\(R_\mathrm{int}^{(\mathrm{SK})}(J_v) = 1 - \alpha = \sigma\,J_v/(B+J_v)\).
	Substituting \eqref{eq:Cp-from-SK} into \eqref{eq:bc-balance}
	gives the corresponding Robin relation
	\begin{equation}
		D_A\,\frac{\partial C_A}{\partial n}\bigg|_{w}
		= \frac{\sigma\,J_v^{\,2}}{B + J_v}\,C_w .
		\label{eq:SK-Robin}
	\end{equation}
\end{itemize}

\paragraph{Implementation as a Robin (mixed) boundary condition.}
Both \eqref{eq:Rint-Robin} and \eqref{eq:SK-Robin} have the canonical
Robin form
\begin{equation}
	k\,(C_w - C_i) \;=\; h_\mathrm{phys}\,C_w ,
	\qquad
	k \equiv D_A/\delta_n ,
	\label{eq:robin-canonical}
\end{equation}
with
\(h_\mathrm{phys} = R_\mathrm{int}\,J_v\) (intrinsic-rejection closure)
or
\(h_\mathrm{phys} = \sigma\,J_v^{\,2}/(B + J_v)\)
(solute-permeability closure). Solving for \(C_w\),
\begin{equation}
	C_w \;=\; \frac{k}{k - h_\mathrm{phys}}\,C_i .
	\label{eq:Cw-closed}
\end{equation}
In OpenFOAM this is realised through the
\texttt{mixedFvPatchScalarField} machinery, which evaluates the wall
value as
\(C_w = \alpha_\mathrm{vF}\,C_\mathrm{ref} + (1-\alpha_\mathrm{vF})\,(C_i + g\,\delta_n)\).
With reference value and reference gradient set to zero, the
implementation reproduces \eqref{eq:Cw-closed} by setting
\begin{equation}
	\alpha_\mathrm{vF} \;=\; -\,\frac{h_\mathrm{phys}}{k - h_\mathrm{phys}} ,
\end{equation}
which is \emph{negative-valued by construction} under physical
operating conditions
(\(0 < h_\mathrm{phys} < k\)). The OpenFOAM identifier
\texttt{valueFraction} therefore stores an algebraic Robin weight
rather than a fraction in \([0,1]\), and negative values are correct.

\paragraph{Three modeling choices per species.}
Equations \eqref{eq:jv-s-k}--\eqref{eq:js-s-k} contain four membrane
parameters: \(A_h\), \(\sigma\), \(B\), and the rejection (linked to
\(\sigma\) and \(B\) via \eqref{eq:Rint-def}). Different experimental
scenarios fix different combinations of these parameters, and
\texttt{centriMembSolver} therefore offers three modeling choices per
species:
\begin{itemize}
	\item Solvent flux (\(\mathbf{U}\) BC, \S\ref{sec:bc-solvent}): the
	Spiegler--Kedem model fixes \(A_h\) and \(\sigma\); the
	constant-flux model fixes \(J_v\) directly; the
	target-average-flux model tunes \(A_h\) on-the-fly to match a
	prescribed \(\langle J_v\rangle\).
	\item Solute flux (\(C_A\) BC, \S\ref{sec:bc-solute}): the
	intrinsic-rejection model fixes \(R_\mathrm{int}\) and uses
	\eqref{eq:Rint-Robin}; the solute-permeability model fixes
	\(B\) (and optionally \(\sigma\)) and uses \eqref{eq:SK-Robin};
	the observed-rejection model tunes \(R_\mathrm{int}\)
	on-the-fly to match a prescribed \(R_\mathrm{obs}\).
	\item Tracer flux (\(C_T\) BC, \S\ref{sec:bc-tracer}): same three
	options applied to a passive tracer with its own constant
	diffusivity \(D_T\) and parameters
	(\(R_{T,\mathrm{int}}\), \(B_T\), \(\sigma_T\)).
\end{itemize}

\subsection{Solvent flux boundary (\texttt{U})}
\label{sec:bc-solvent}
The solvent flux condition determines the normal component of
\(\mathbf{U}\) at the membrane.  Three models are provided.  In all
three the local solvent flux is clipped to non-negative values,
\(J_v \leftarrow \max(0, J_v)\), so that the membrane never admits
permeate-to-feed back-flow during transient or extreme conditions;
this safeguard never affects steady-state results in the validated
operating envelope but stabilises early-time behaviour while the
dynamic pressure has not yet established.

\paragraph{Spiegler–Kedem model.}
Implements Eq.~\eqref{eq:jv-s-k}.  Dictionary entries include the
hydraulic permeability \texttt{Ah}, the reflection coefficient
\texttt{sigma}, virial coefficients \texttt{virialCoeffs} for the
osmotic pressure (with \(C_A\) in \si{kg/m^3} and \(\pi\) in
\si{Pa}), and an optional permeate-side pressure reference
\texttt{pPermConst}.  The example below sets
\texttt{Ah = 0.5e-11}, \texttt{sigma = 1.0} and a virial expansion
of the form \(\pi(C_A) = a_1 C_A + a_2 C_A^2\):

\begin{verbatim}
membraneA
{
    type    membraneSolventFlux;
    model   SpieglerKedem;
    Ah      0.5e-11;        // [m/(Pa s)]
    sigma   1.0;            // [-]
    virialCoeffs (2.5e3 0); // a1 [Pa m3 kg-1], a2 [Pa m6 kg-2]
    pPermConst 0;           // [Pa]
    value  uniform (0 0 0);
}
\end{verbatim}

\paragraph{Constant flux model.}
Prescribes a uniform solvent flux \(J_v\) independent of driving
forces.  Useful for benchmarking and when the membrane permeability is
unknown.  Example:
\begin{verbatim}
membraneA
{
    type    membraneSolventFlux;
    model   constantFlux;
    Jv      3.0e-5;        // [m/s]
    value   uniform (0 0 0);
}
\end{verbatim}

\paragraph{Target average flux model.}
Tunes an effective permeability \(A_h^*\) such that the area-averaged
flux matches a target \texttt{JvAverage}.  Driving forces are still
proportional to \(\Delta p - \sigma\Delta\pi\) as in Spiegler--Kedem,
but \(A_h^*\) is updated each time step within user-specified bounds
\texttt{AhMin} and \texttt{AhMax}.  Particularly useful for calibrating
models against a measured average permeate flux.  Example:
\begin{verbatim}
membraneA
{
    type        membraneSolventFlux;
    model       targetAverageFlux;
    JvAverage   3.0e-5;        // target <Jv> [m/s]
    AhMin       1.0e-13;       // lower bound on A_h^* [m/(Pa s)]
    AhMax       1.0e-9;        // upper bound on A_h^*
    sigma       1.0;
    virialCoeffs (7314.5 8.253);
    pPermConst  0;
    value       uniform (0 0 0);
}
\end{verbatim}

\paragraph{Diagnostics.}
For each membrane patch the solver writes the instantaneous and
time–windowed area–average of \(J_v\).  If window averaging is
enabled (\texttt{windowAverage true}), the time scale \texttt{Twindow}
sets the averaging duration.  The internal field \texttt{Jv} is
available for further post–processing.

\subsection{Solute flux boundary (\texttt{CA})}
\label{sec:bc-solute}
The solute boundary condition couples the solute transport to the
solvent flux through the local mass balance derived in
\S\ref{sec:bc-theory} (Eqs.~\eqref{eq:bc-balance} and following), and
is implemented as an OpenFOAM \texttt{mixedFvPatchScalarField} (Robin
BC). Three modeling closures are available.

\paragraph{Intrinsic rejection model.}
Fixes \(R_\mathrm{int}\) directly; uses Eqs.~\eqref{eq:Cp-from-Rint}
and~\eqref{eq:Rint-Robin}. Dictionary entries: \texttt{model
intrinsicRejection}, \texttt{Rint}, and the reference bulk
concentration \texttt{CAb} used for the polarisation diagnostic
\(\Gamma\).

\paragraph{Solute permeability model.}
Fixes the solute permeability \(B\) (and optionally \(\sigma\)); uses
Eqs.~\eqref{eq:Cp-from-SK} and~\eqref{eq:SK-Robin}. Reduces to a
solution--diffusion model when \(\sigma = 0\). If \texttt{sigma} is
specified it overrides the reflection coefficient from the solvent
BC; otherwise \(\sigma\) is inherited from the solvent BC on
\texttt{U}, defaulting to \(1\) if no membrane solvent BC is found.
Dictionary entries: \texttt{model solutePermeability}, \texttt{B},
optionally \texttt{sigma}, and \texttt{CAb}.

\paragraph{Observed rejection model.}
Targets a prescribed observed rejection
\begin{equation}
	R_\mathrm{obs} \;=\; 1 - \frac{\langle J_v C_p \rangle}{\langle J_v \rangle\,C_b} ,
	\label{eq:Robs-def}
\end{equation}
where \(\langle\cdot\rangle\) denotes an area average and \(C_b\) is
the bulk reference concentration. The solver iteratively determines
an intrinsic rejection \(R_\mathrm{int}\) such that
\(R_\mathrm{obs}\) matches the target within a specified tolerance,
applying Eq.~\eqref{eq:Rint-Robin} face-by-face with the inferred
\(R_\mathrm{int}\). Dictionary entries: \texttt{model
observedRejection}, \texttt{Robs} (target), \texttt{RintMin},
\texttt{RintMax}, \texttt{tol}, \texttt{maxIter}, and \texttt{CAb}.
Useful when matching experimental rejection measurements.

\paragraph{Diagnostics.}
Per-face fluxes \(J_s\), \(J_v\) and permeate concentration \(C_p\)
are exposed via the patch fields \texttt{Js}, \texttt{Jv}, and
\texttt{CAp}, and logged in
\texttt{postProcessing/membraneSoluteFlux/<patch>/Rint\_vs\_time.dat}
(columns: \emph{time}, \(R_\mathrm{int}\), \(R_\mathrm{obs}\)). The
solver also writes the local concentration polarisation factor
\(\Gamma = (C_w - C_b)/C_b\) to the volume field \texttt{Gama}.


\subsection{Tracer flux boundary (\texttt{CT})}
\label{sec:bc-tracer}
The tracer boundary condition reuses the same Robin BC machinery
described in \S\ref{sec:bc-theory} (Eqs.~\eqref{eq:bc-balance}--\eqref{eq:Cw-closed}),
applied to the passive tracer field \(C_T\) with diffusivity \(D_T\)
and tracer-specific parameters \(R_{T,\mathrm{int}}\), \(B_T\),
\(\sigma_T\). All operations on \(C_T\) are independent of those on
\(C_A\); the only coupling is through the shared solvent flux
\(J_v\). Three modeling closures are available, in direct analogy to
the solute case.

\paragraph{Intrinsic rejection model.}
Tracer counterpart of \eqref{eq:Cp-from-Rint} and
\eqref{eq:Rint-Robin}: \(C_{T,p} = (1 - R_{T,\mathrm{int}})\,C_{T,w}\)
and \(D_T\,\partial C_T/\partial n|_{w} = R_{T,\mathrm{int}}\,J_v\,C_{T,w}\).
Dictionary entries: \texttt{model intrinsicRejection},
\texttt{RTint}, \texttt{CTb}.

\paragraph{Solute permeability model.}
Tracer counterpart of the SK closure
\eqref{eq:Cp-from-SK}--\eqref{eq:SK-Robin} with \((\sigma, B)\)
replaced by \((\sigma_T, B_T)\). When \texttt{sigmaT} is omitted, the
default is \(\sigma_T = 1\) (full rejection of the tracer); unlike
the solute case, \(\sigma_T\) is \emph{not} inherited from the
solvent BC, since the reflection coefficient of a passive tracer is
generally unrelated to that of the primary solute. Dictionary
entries: \texttt{model solutePermeability}, \texttt{BT}, optionally
\texttt{sigmaT}, \texttt{CTb}.

\paragraph{Observed rejection model.}
Tracer counterpart of \eqref{eq:Robs-def}, with \((C_p, C_b)\)
replaced by \((C_{T,p}, C_{T,b})\). The solver iteratively determines
\(R_{T,\mathrm{int}}\) such that the achieved
\(R_{T,\mathrm{obs}}\) matches a prescribed target. Dictionary
entries: \texttt{model observedRejection}, \texttt{RTobs},
\texttt{RTintMin}, \texttt{RTintMax}, \texttt{tol}, \texttt{maxIter},
\texttt{CTb}.

\paragraph{Tracer dictionary example.}
\begin{verbatim}
membraneB
{
    type     membraneTracerFlux;
    model    intrinsicRejection;
    RTint    1.0;          // perfectly rejected tracer
    CTb      1.0e-3;       // reference [kg/m^3]
    value    uniform 0;
}
\end{verbatim}

\paragraph{Diagnostics.}
Per-face tracer flux \(J_{s,T}\), solvent flux \(J_v\) and permeate tracer concentration \(C_{T,p}\) are exposed via the patch fields \texttt{JsT}, \texttt{Jv} and \texttt{CTp}.  The solver writes the observed tracer rejection and concentration polarisation factor \(\Gamma_T = (C_{T,w}-C_{T,b})/C_{T,b}\) to the summary files for post-processing.

\subsection{Membrane patch dictionaries}


Tables~\ref{tab:membrane-solvent} and~\ref{tab:membrane-solute}
summarise the principal dictionary entries for solvent and solute
boundary conditions.  Most values default to physically reasonable
numbers; only those shown as required must be provided.

\begin{table}[h]
  \centering
  \caption{Summary of membrane solvent flux boundary entries.}
  \label{tab:membrane-solvent}
  \begin{tabular}{@{}llp{8cm}@{}}
    \toprule
    Key & Units & Description\\
    \midrule
    \texttt{model} & – & \texttt{SpieglerKedem}, \texttt{constantFlux} or \texttt{targetAverageFlux}\\
    \texttt{Ah} & \si{m/(Pa\,s)} & Hydraulic permeability for Spiegler–Kedem\\
    \texttt{sigma} & – & Reflection coefficient \(\sigma\); defaults to 1\\
    \texttt{virialCoeffs} & \si{Pa\,m^{3k}\,kg^{-k}} & List \((a_1,a_2,\dots)\) for osmotic pressure $\pi(C_A)=\sum_k a_k C_A^k$ ($a_1$ in \si{Pa\,m^3\,kg^{-1}}, $a_2$ in \si{Pa\,m^6\,kg^{-2}}, etc.)\\
    \texttt{Jv} & \si{m/s} & Constant flux magnitude (constantFlux model)\\
    \texttt{JvAverage} & \si{m/s} & Target area–average flux (targetAverageFlux)\\
    \texttt{AhMin}, \texttt{AhMax} & \si{m/(Pa\,s)} & Limits on tuned permeability\\
    \texttt{pPermConst} & \si{Pa} & Permeate-side reference pressure (default 0)\\
    \texttt{windowAverage}, \texttt{Twindow} & – & Enable and set window for time–averaged output\\
    \bottomrule
  \end{tabular}
\end{table}

\begin{table}[h]
  \centering
  \caption{Summary of membrane solute flux boundary entries.}
  \label{tab:membrane-solute}
  \begin{tabular}{@{}llp{8cm}@{}}
    \toprule
    Key & Units & Description\\
    \midrule
    \texttt{model} & – & \texttt{intrinsicRejection}, \texttt{solutePermeability} or \texttt{observedRejection}\\
    \texttt{Rint} & – & Intrinsic rejection coefficient (intrinsicRejection)\\
    \texttt{B} & \si{m/s} & Solute permeability (solutePermeability)\\
    \texttt{sigma} & – & Optional reflection coefficient (overrides solvent BC; default 1)\\
    \texttt{Robs} & – & Target observed rejection (observedRejection)\\
    \texttt{CAb} & \si{kg/m^3} & Reference bulk concentration for diagnostics\\
    \texttt{RintMin}, \texttt{RintMax} & – & Bounds for inverse rejection search\\
    \texttt{tol}, \texttt{maxIter} & – & Tolerance and maximum iterations for secant/bisection\\
    \texttt{UName} & – & Name of velocity field (defaults to \texttt{U})\\
    \bottomrule
  \end{tabular}
\end{table}

\begin{table}[h]
  \centering
  \caption{Summary of membrane tracer flux boundary entries.}
  \label{tab:membrane-tracer}
  \begin{tabular}{@{}llp{8cm}@{}}
    \toprule
    Key & Units & Description\\
    \midrule
    \texttt{model} & -- & \texttt{intrinsicRejection}, \texttt{solutePermeability} or \texttt{observedRejection}\\
    \texttt{RTint} & -- & Intrinsic tracer rejection coefficient (intrinsicRejection)\\
    \texttt{BT} & \si{m/s} & Tracer permeability (solutePermeability)\\
    \texttt{sigmaT} & -- & Optional tracer reflection coefficient (default 1; full rejection)\\
    \texttt{RTobs} & -- & Target observed tracer rejection (observedRejection)\\
    \texttt{CTb} & \si{kg/m^3} & Reference bulk tracer concentration for diagnostics\\
    \texttt{RTintMin}, \texttt{RTintMax} & -- & Bounds for inverse rejection search\\
    \texttt{tol}, \texttt{maxIter} & -- & Tolerance and maximum iterations for secant/bisection\\
    \texttt{UName} & -- & Name of velocity field (defaults to \texttt{U})\\
    \bottomrule
  \end{tabular}
\end{table}


% ===========================================================================


\section{Solver structure and implementation}
\label{sec:solver-structure}

The \texttt{centriMembSolver} follows a PIMPLE--type pressure--velocity
coupling, extended with (i) concentration--dependent transport
properties, (ii) a rotating reference frame, and (iii) custom membrane
boundary conditions and diagnostics. The overall time loop is driven by
\texttt{runTime.run()}, and at each time step a nested PIMPLE loop
solves momentum, pressure and solute transport. Additional headers
perform pressure reconstruction and conservation monitoring before the
fields are written.

\subsection{Overall time loop and custom PIMPLE algorithm}

For readers familiar with C++, the main structure of
\texttt{centriMembSolver.C} can be summarised schematically as

\begin{verbatim}
	createTime;
	createMesh;
	createFields;               // U, p_rgh, CA, Js, Jv, CAp, Gama,
	                            // and tracer fields (CT, DT, JsT,
	                            // CTp, GamaT) if solveTracer=true
	createPimpleControl;        // pimpleControl
	createRDeltaT;              // LTS auto-detection (sets bool LTS)
	updateTransportProperties;  // nu(CA), DAB(CA), rhok

	while (runTime.run())
	{
		// LTS or CFL-based time-step control
		if (LTS) { #include "setRDeltaT.H" }
		else     { #include "CourantNo.H"; #include "setDeltaT.H" }

		++runTime;
		Info<< "Time = " << runTime.timeName() << endl;

		while (pimple.loop())           // outer PIMPLE loop
		{
			#include "UEqn.H"           // assemble & solve momentum
			#include "CAEqn.H"          // assemble & solve solute
			#include "updateTransportProperties.H"

			while (pimple.correct())    // inner pressure loop
			{
				#include "pEqn.H"
			}

			if (pimple.turbCorr())
			{
				turbulence->correct();
			}
		}

		#include "CTEqn.H"
		#include "pReconstruct.H"
		#include "monitorConservation.H"
		#include "Gama.H"
		#include "GamaT.H"

		runTime.write();
		runTime.printExecutionTime(Info);
	}
\end{verbatim}

Compared with a ``standard'' OpenFOAM transient solver:

\begin{itemize}
	\item momentum \texttt{(UEqn.H)} and solute \texttt{(CAEqn.H)} are
	solved together at the beginning of each PIMPLE outer iteration;
	\item transport properties (viscosity, diffusivity, density) are
	updated \emph{inside} the PIMPLE loop via
	\texttt{updateTransportProperties.H}, so they respond
	iteratively to the current concentration field;
	\item several post--processing steps (absolute pressure, conservation,
	membrane diagnostic \(\Gamma\), body forces) are executed at the
	end of every time step.
\end{itemize}

This design makes the coupling between flow, solute transport and
membrane behaviour explicit while still using the familiar PIMPLE
machinery.

\subsection{Field creation, membrane diagnostics and transport properties}

Field creation is handled in two main headers:

\begin{itemize}
	\item \textbf{\texttt{createFields.H}} reads the primary fields
	(\texttt{U}, \texttt{p\_rgh}, \texttt{CA}, etc.) from the
	\texttt{0/} directory and the physical parameters from
	\texttt{constant/transportProperties} and
	\texttt{constant/SRFProperties}. It also constructs the
	turbulence model and creates all membrane diagnostic fields
	directly: \texttt{Js} (solute flux), \texttt{Jv} (solvent flux),
	\texttt{CAp} (permeate solute concentration) and \texttt{Gama}
	(dimensionless polarisation factor). When the user enables the
	tracer via the runtime switch \texttt{solveTracer} in
	\texttt{controlDict}, \texttt{createFields.H} additionally
	allocates the tracer fields \texttt{CT}, \texttt{DT},
	\texttt{DTeff}, \texttt{JsT}, \texttt{CTp} and \texttt{GamaT}
	through \texttt{autoPtr} holders. All diagnostic fields are
	created with \texttt{AUTO\_WRITE} so that they are automatically
	written at output times.
\end{itemize}

The header \texttt{updateTransportProperties.H} is called once before
the time loop and then at each outer PIMPLE iteration. It performs three
updates:

\begin{enumerate}
	\item \emph{Viscosity}: the laminar transport model
	\texttt{laminarTransport} is corrected using the current
	concentration field \texttt{CA}, so that the kinematic viscosity
	\(\nu(\mathrm{CA})\) follows the chosen \texttt{CAViscosity}
	model.
	
	\item \emph{Molecular diffusivity}: the scalar field \texttt{DAB} is
	updated as a function of \texttt{CA}, using a polynomial
	relation. Boundary conditions on \texttt{DAB} are corrected
	afterwards.
	
	\item \emph{Density}: a non--dimensional density field \texttt{rhok}
	is recomputed according to a Boussinesq--type law,
	\(\rho_k = 1 + \beta\,\mathrm{CA}\), and its boundary
	conditions are corrected.
\end{enumerate}

Updating these properties inside the PIMPLE loop strengthens the
coupling between momentum and solute transport: any change in
\texttt{CA} is immediately reflected in \(\nu\), \(D_{AB}\) and
\(\rho_k\).

\subsection{Momentum and pressure equations}

Inside the outer \texttt{while (pimple.loop())}:

\begin{itemize}
	\item \textbf{\texttt{UEqn.H}} assembles the transient momentum
	equation using implicit time discretisation
	(\texttt{fvm::ddt}), a convection term
	\texttt{fvm::div(phi,U)}, viscous/turbulent stresses via
	\texttt{turbulence->divDevReff(U)}, and a Coriolis source term
	\(2\rho\,\boldsymbol{\Omega}\times\mathbf{U}\) for the rotating
	frame. The equation is then solved for the velocity field
	\texttt{U}.
	
	\item \textbf{\texttt{CAEqn.H}} assembles and solves the unsteady
	advection--diffusion equation for the solute concentration
	\texttt{CA}. The temporal and spatial discretisations follow
	the user-specified schemes in \texttt{system/fvSchemes}; for the
	tutorials a bounded convection scheme is recommended together
	with the time scheme appropriate for the chosen ddt
	(\texttt{Euler}, \texttt{CrankNicolson}, or \texttt{localEuler}
	when LTS is enabled).
\end{itemize}

After these solves, \texttt{updateTransportProperties.H} is called again
to refresh \(\nu(\mathrm{CA})\), \(D_{AB}(\mathrm{CA})\) and \(\rho_k\)
based on the newly updated \texttt{CA} field.

The pressure equation is then handled in the inner loop
\texttt{while (pimple.correct())} through \texttt{pEqn.H}. This step
enforces continuity, updates \texttt{p\_rgh} (and the associated face
flux \texttt{phi}) and accounts for the rotation--induced hydrostatic
term and the specified pressure reference. When
\texttt{pimple.turbCorr()} is true, the turbulence model is also
corrected in the same outer iteration.

\subsection{Membrane boundary conditions and the \texorpdfstring{$\Gamma$}{Gamma} diagnostic}

The actual membrane physics is implemented in custom boundary condition
classes applied to \texttt{U} and \texttt{CA}. These boundary
conditions:

\begin{itemize}
	\item compute, for each membrane face, the solvent flux \(J_v\) and
	solute flux \(J_s\) according to the selected transport model
	(e.g. pressure--driven with osmotic correction);
	\item update the permeate concentration \(C_p\) and a dimensionless
	diagnostic \(\Gamma\) that characterises local solute rejection
	or concentration polarisation.
\end{itemize}

At volume level, these per--face quantities are exported into the
membrane diagnostic fields created in \texttt{createFields.H}
(\texttt{Js}, \texttt{Jv}, \texttt{CAp}, \texttt{Gama}, and the
optional tracer fields when \texttt{solveTracer=true}). The header
\texttt{Gama.H} is responsible for copying the per--face \(\Gamma\)
values from the membrane solute boundary condition into the boundary of
the volume field \texttt{Gama}, so that \texttt{Gama} can be visualised
and sampled like any other \texttt{volScalarField}. For efficiency,
\texttt{Gama.H} caches the membrane patch indices and a pointer to the
\texttt{CA} field the first time it is called.

\subsection{Post-processing: pressure reconstruction, conservation and diagnostics}

After the PIMPLE loop converges at a given time, the following
post--processing headers are executed before writing:

\paragraph{\texttt{pReconstruct.H}}

The solver works internally with the modified pressure \texttt{p\_rgh}
to simplify the treatment of gravity in a rotating frame. For analysis,
it is convenient to have the absolute pressure field. The header
\texttt{pReconstruct.H} reconstructs

\[
p = (p\_\text{rgh} + \rho_k\,g\,h)\,\rho_0 + p_\text{ref},
\]

where \(\rho_0\) is a reference density and \(p_\text{ref}\) the chosen
pressure reference, and then calls \texttt{p.correctBoundaryConditions()}
to keep the boundaries consistent.

\paragraph{\texttt{monitorConservation.H}}

This header checks global solute conservation in a parallel--safe way.
It requires the fields \texttt{CA}, \texttt{U}, \texttt{Js},
\texttt{Jv}, \texttt{CAp} and the face flux \texttt{phi}; an optional
\texttt{Deff} field can be supplied to account for diffusive fluxes at
walls. It builds (or appends to) the file

\[
\texttt{postProcessing/conservation/cons.dat},
\]

writing at each output time the total solute mass in the domain, the
time derivative of that mass, the integrated inlet, outlet, membrane and
diffusive contributions, and a residual and relative error that measure
how well the discrete balance is satisfied. This file is the main tool
to verify that the coupled flow--membrane scheme is globally
conservative.

\paragraph{\texttt{Gama.H} and \texttt{GamaT.H}}

As described above, \texttt{Gama.H} maps the per--face diagnostic
\(\Gamma\) from the membrane solute boundary condition into the volume field
\texttt{Gama}.  In the 2026 release, a companion header \texttt{GamaT.H}
performs the analogous mapping for the tracer polarisation factor \texttt{GamaT}.
This is done patch by patch on the membrane boundaries, with simple consistency checks (same number of faces) to avoid copying errors.

At the end of this post--processing stage,
\texttt{runTime.write()} writes all \texttt{AUTO\_WRITE} fields
(\texttt{U}, \texttt{p}, \texttt{p\_rgh}, \texttt{CA}, \texttt{Js},
\texttt{Jv}, \texttt{CAp}, \texttt{Gama}, and the tracer
counterparts when enabled), and the solver advances to the next
time level.

% ===========================================================================


\section{Case studies with \texttt{centriMembSolver}}
\label{sec:cases}

This section presents an example case that illustrates how to use
\texttt{centriMembSolver}.  It is designed to be readable even for
readers who are not familiar with OpenFOAM.

In OpenFOAM, each simulation (``case'') is stored in a directory with a
standard structure:
\begin{itemize}
	\item \texttt{0/} contains the \emph{initial and boundary conditions}
	for all fields (for example velocity, pressure and concentration);
	\item \texttt{constant/} contains the \emph{physical properties}
	(for example fluid properties, solute properties, membrane
	parameters, rotation settings);
	\item \texttt{system/} contains the \emph{numerical setup}
	(mesh generation with \texttt{blockMesh}, time step controls,
	linear solvers, etc.).
\end{itemize}

To run a case, one typically:
\begin{enumerate}
	\item generates the mesh (for example with \texttt{blockMesh} or a
	helper script);
	\item runs the solver (here \texttt{centriMembSolver});
	\item inspects the results in ParaView or in text files written in
	\texttt{postProcessing/}.
\end{enumerate}

Below we describe a simple 2D slit case (\texttt{rotatingSlit})
that exercises the full pipeline (mesh, BCs, solver, post-processing).

% ---------------------------------------------------------------------------
\subsection{\texttt{rotatingSlit}}

The example case \texttt{rotatingSlit} distributed with the solver
provides a simple yet instructive demonstration.  It consists of a
two-dimensional slit domain rotated about its centre.  The membrane
occupies a section of the side wall; the remaining boundaries are
inlet, outlet, symmetry and empty patches.

\subsubsection*{Geometry and mesh}

The geometry is defined by \texttt{system/blockMeshDict}.  A narrow
slit of length \SI{4}{\centi\metre} and width \SI{1}{\milli\metre} is
discretised into a structured mesh.  One boundary, named
\texttt{membraneB}, is assigned the membrane boundary conditions.  A
second boundary \texttt{membraneA} is present but disabled in this
case.  Front and back patches are declared as \texttt{empty} to
produce a two-dimensional simulation.

\subsubsection*{Initial and boundary conditions}

The inlet velocity is fixed at $\mathbf{U}_\mathrm{in} = (0.01, 0, 0)$~m/s
along the slit axis.  The initial solute
concentration is uniform (\SI{5}{kg/m^3}) and the inlet condition sets
\texttt{CAin} to a higher or lower value to induce concentration
polarisation.  The membrane on patch \texttt{membraneB} is configured
with
\begin{itemize}
	\item a target average solvent flux of \SI{3e-5}{m/s},
	\item limits \texttt{AhMin = 1e-13} and \texttt{AhMax = 1e-9},
	\item reflection coefficient $\sigma=1$, and
	\item virial coefficients
	$(7314.5\,\mathrm{Pa\,m^3\,kg^{-1}},
	8.253\,\mathrm{Pa\,m^6\,kg^{-2}})$.
\end{itemize}
On the solute side the membrane uses the intrinsic rejection model with
\texttt{Rint = 0.99} and reference concentration \texttt{CAb = 1.0}.  A
fixed–flux pressure condition ensures continuity of volumetric flow.

\subsubsection*{Running the case}

To run the example, follow these steps in the case directory:

\begin{enumerate}
	\item Compile the solver (only needed once) by executing \texttt{wmake}
	in the \texttt{solver} directory.
	\item Generate the mesh using
	\begin{verbatim}
		blockMesh
	\end{verbatim}
	\item Adjust \texttt{constant/SRFProperties} to set the rotation
	speed $\boldsymbol{\Omega}$ if desired.  The default example uses
	\texttt{rpm = 0}, but any non–zero value can be specified.
	\item Run the simulation with
	\begin{verbatim}
		centriMembSolver -case rotatingSlit > log &
	\end{verbatim}
	\item Monitor progress via the log or the files in
	\texttt{postProcessing}.  Key outputs include
	\texttt{Rint\_vs\_time.dat} (intrinsic rejection over time),
	\texttt{calibratedProperties.dat} (effective permeability for
	target average flux) and \texttt{cons.dat} (mass conservation).
\end{enumerate}

Figure~\ref{fig:rotatingSlit-results} shows a real run of the bundled
\texttt{rotatingSlit} tutorial.  Two transients are visible.

\textbf{Panel (a) --- targetAverageFlux calibration.}
The \texttt{targetAverageFlux} solvent model dynamically adjusts the
effective hydraulic permeability $A_h^\ast$ at every time step so that
the area-averaged solvent flux matches the prescribed
\texttt{JvAverage}.  Starting from a low initial value (no
polarisation, weak osmotic counter-pressure), $A_h^\ast$ converges to
its steady operating value (here $4.36\times 10^{-11}$~m\,(Pa\,s)$^{-1}$)
within ${\sim}15$~s of simulated time.  By construction the model
keeps $\langle J_v\rangle$ exactly at the target throughout.

\textbf{Panel (b) --- concentration polarisation signature.}
The intrinsic rejection imposed at the membrane is constant by
construction ($R_\mathrm{int} = 0.99$, dashed line), while the
observed rejection
\begin{equation*}
	R_\mathrm{obs}(t) = 1 - \frac{\langle J_v C_p\rangle(t)}{\langle J_v\rangle(t)\,C_b}
\end{equation*}
descends from values close to $R_\mathrm{int}$ at the onset of
permeation (no polarisation yet) towards a steady value
substantially below $R_\mathrm{int}$ ($R_\mathrm{obs}\approx 0.60$
here), reflecting the equilibrium polarisation that the simulation
establishes at the membrane.  The $R_\mathrm{int}{-}R_\mathrm{obs}$
gap is the operational signature of concentration polarisation, and
is automatically logged under
\texttt{postProcessing/membraneSoluteFlux/<patch>/Rint\_vs\_time.dat}
(columns: time, $R_\mathrm{int}$, $R_\mathrm{obs}$).

\begin{figure}[h]
	\centering
	\includegraphics[width=\linewidth]{figures/rotatingSlit_results.pdf}
	\caption{Real outputs of the bundled \texttt{rotatingSlit} tutorial
		(2D slit, $L\times h = $ \SI{4}{cm}$\times$\SI{1}{mm}, NaCl-like
		feed at \SI{5}{kg/m^3}, $\sigma=1$, $R_\mathrm{int}=0.99$,
		\texttt{JvAverage}$= 3\times 10^{-5}$~m/s).
		(a) Calibrated hydraulic permeability $A_h^\ast(t)$ extracted
		from the patch dictionary at each write time, showing the
		\texttt{targetAverageFlux} model converging to the steady
		operating permeability ($\sim 4.36\times 10^{-11}$~m\,(Pa\,s)$^{-1}$).
		(b) Time evolution of the observed rejection $R_\mathrm{obs}$
		on patch \texttt{membraneB} (red), starting near
		$R_\mathrm{int}$ (dashed line, prescribed) and stabilising
		around $0.60$ as the polarised layer builds up.  Data from
		\texttt{postProcessing/membraneSoluteFlux/membraneB/Rint\_vs\_time.dat}
		and the saved patch dictionaries (\texttt{<time>/U}).}
	\label{fig:rotatingSlit-results}
\end{figure}

% ===========================================================================
\section{Practical guidance}

\paragraph{Stability and time stepping.}
With the Coriolis source treated explicitly, the convective Courant
condition is the dominant time-step constraint; keep the Courant
number below \(\sim 1\) near the membrane.  When using
\texttt{targetAverageFlux}, choose the bounds \texttt{AhMin} and
\texttt{AhMax} narrowly around the expected steady operating
permeability: a range that is too wide allows the calibrated
\(A_h^*\) to wander outside the physical regime and may produce
oscillatory behaviour.  When using \texttt{observedRejection}, pick a
tight tolerance \texttt{tol} and verify that the target
\(R_\mathrm{obs}\) lies inside the physically realistic interval
\([\,R_{\rm int,min},\,R_{\rm int,max}\,]\) -- otherwise the inverse
search hits the bracket and reports a clamped value rather than the
target.

\paragraph{Parallel consistency.}
All membrane boundary conditions and diagnostics use parallel
reductions (\texttt{Pstream::reduce}, \texttt{returnReduce}) and are
safe out of the box on decomposed meshes.  For complex geometries or
very fine meshes, however, it is advisable to invoke
\texttt{redistributePar} after \texttt{decomposePar} so that each
membrane patch is owned by a single processor: this avoids subtle
load-balancing artefacts in the area-weighted averages.  The
conservation monitor (\texttt{postProcessing/conservation/cons.dat})
provides a quick sanity check on mass balances across processor
boundaries.

% ===========================================================================
\section{References}

\begin{enumerate}
	\item Spiegler, K. S. \& Kedem, O. (1966). Thermodynamics of hyperfiltration (reverse osmosis): criteria for efficient membranes. \emph{Desalination}, 1, 311–326.
	\item Wu, J. J. (2019). On the application of the Spiegler–Kedem model to forward osmosis. \emph{BMC Chemical Engineering}, 1, article 15. \href{https://doi.org/10.1186/s42480-019-0014-8}{doi:10.1186/s42480-019-0014-8}. (Discusses how the S–K model accommodates deviations from ideal semi-permeability.)
	\item Chandrasekhar, S. (1961). \emph{Hydrodynamic and Hydromagnetic Stability}. Clarendon Press, Oxford. International Series of Monographs on Physics. (Classical reference for absorbing the centrifugal term into the pressure gradient in rotating Boussinesq flows.)
	\item Lopez, J. M., Marques, F. \& Avila, M. (2013). The Boussinesq approximation in rapidly rotating flows. \emph{Journal of Fluid Mechanics}, 737, 56–77. \href{https://doi.org/10.1017/jfm.2013.558}{doi:10.1017/jfm.2013.558}. (Explicit modern formulation of the gravity + centrifugal absorption; eqs.~(2.1)–(2.3) used as reference for the \texttt{gh}/\texttt{p\_rgh} construction in this solver.)
\end{enumerate}

\end{document}
